\subsection{Linear Prediction and Generalized Linear Models}

Linear prediction is the local model against which many nonlinear learning methods are compared.

\begin{definition} \label{def:sl-linear-predictor} \hlt{Linear Predictor}\\
For features $x\in\R^p$, a linear score has the form
\[
f_{w,b}(x)=w^\top x+b.
\]
Regression interprets this score directly, while classification often passes it through a sign, sigmoid, or softmax map.
\end{definition}

\begin{example} \label{ex:sl-linear-regression-risk} \hlt{Least-Squares Linear Regression}\\
Given design matrix $X$ and response vector $y$, ordinary least squares solves
\[
\widehat{\beta}\in\operatorname*{arg\,min}_{\beta}\norm{y-X\beta}_2^2.
\]
When $X^\top X$ is invertible, $\widehat{\beta}=(X^\top X)^{-1}X^\top y$.
\end{example}

\begin{definition} \label{def:sl-generalized-linear-model} \hlt{Generalized Linear Model}\\
A generalized linear model specifies a conditional distribution from an exponential family and links its mean to a linear score:
\[
g(\mathbb{E}[Y\mid X=x])=w^\top x+b.
\]
Logistic regression is the Bernoulli case with logit link.
\end{definition}

\begin{remark} \label{rem:sl-feature-maps-linear-models} \hlt{Feature Maps}\\
Applying a linear model to transformed features $\phi(x)$ gives nonlinear prediction in the original input space.  Kernel methods and neural networks can both be viewed as systematic ways to choose or learn such feature maps.
\end{remark}
